% Evidence-backed Phase 1 section. Integrate into the official seven-phase
% Report.tex template, retaining its compulsory top-level section order.
% This section was drafted with Codex assistance; disclose that truthfully.

\subsection{Phase 1 foundation: Improved A / LegalSpan}
Phase~1 required a legal language model trained from scratch on the four
provided sources---Acts, High Court judgments, Supreme Court judgments, and
Tribunal documents---to generate a 128-word continuation for each 512-word
prompt. BERTScore F1 was the primary metric; ROUGE-L F1 and BLEU-4 were
secondary. The course supplied no public development split, so the archived
pipeline constructed seeded train, development, and confirmation partitions.
It grouped exact and near-duplicate documents by normalized SHA-256 and
five-word-shingle SimHash before assignment. The official complete corpus was
described as exceeding 4.5 billion tokens; the selected run trained on
209,927,459 supervised continuation tokens, not the full corpus.

The tokenizer was a handwritten, reversible byte-level BPE learner and codec,
fitted on training documents only. It sampled 25 million characters from each
of the four sources and produced a vocabulary of 32,768 tokens. The prepared
stream contained 189,729 training examples, 144 development prompts, and 127
confirmation prompts. It alternated exact 512-word/128-word continuation
windows with nonoverlapping token windows whose prefix was cut at a seeded
position. One quarter of documents used the model's 8,192-token maximum for
token windows; the remainder used 2,047-token windows. Only suffix positions
contributed to the training loss.

LegalSpan contained 112,095,887 parameters. Its backbone was a 12-block
causal Transformer of width 768 with 12 attention heads of width 64,
2,048-dimensional SwiGLU feed-forward layers, pre-RMSNorm, RoPE, and tied
input/output embeddings. The last four blocks added a 128-dimensional
prompt-memory route. At each generation position, a projected query scored
eligible observed prompt tokens, with a learned bias from heuristic prompt
roles and a continuity feature based on up to four already-observed tokens.
The resulting prompt context passed through a sigmoid-gated residual fusion.
Prompt states preceding the continuation boundary remained unchanged. A
second learned gate mixed the normal vocabulary probability with copy mass
aggregated over prompt token IDs:
\begin{equation}
P(y_t\mid x_{\leq t})=(1-g_t)P_{\mathrm{vocab}}(y_t\mid x_{\leq t})
 +g_tP_{\mathrm{copy}}(y_t\mid x_{\leq t}).
\end{equation}
The role cues were surface heuristics, not gold entity labels. The architecture
combined established causal-Transformer and pointer-generator ideas in a
course-specific route/copy hypothesis; it does not claim to invent copying.

The recovered training record specifies AdamW, peak learning rate
$3\times10^{-4}$, betas $(0.9,0.95)$, weight decay $0.1$, 2\% linear warm-up,
cosine decay to 10\% of the peak rate, BF16 autocast, gradient clipping at
$1.0$, and 32 gradient-accumulation examples per optimizer step. The run
completed 11,888 steps and processed 581,555,637 input tokens in about 6.11
hours. Input tokens and supervised target tokens count different quantities.
The pipeline compared checkpoints and three decoding policies on development
data, selecting nucleus sampling with temperature $0.8$, top-$p=0.9$, and
repetition penalty $1.05$. A separate confirmation set was evaluated after
selection. The final generator enforced exactly 128 complete words and wrote
one root-level \texttt{prediction.json} for 1,000 official test IDs.

\begin{table}[H]
\centering
\caption{Historical Phase~1 Codabench scores from preserved screenshots,
not a fresh leaderboard retrieval.}
\begin{tabular}{llc}
\toprule
Lineage & Computation graph & BERTScore F1 \\
\midrule
A0 & Dense causal decoder & 0.8304 \\
B0 & Encoder--decoder Transformer & 0.7972 \\
C & Prompt encoder plus GRU generator & 0.8237 \\
\textbf{Improved A / LegalSpan} & Routed/copy causal decoder & \textbf{0.8341} \\
Improved B & B0 graph with continued task-aligned training & 0.8244 \\
\bottomrule
\end{tabular}
\end{table}
LegalSpan was the highest evidenced Phase~1 submission, 0.0037 above A0.
B1 and C1 own-BPE reruns had valid prediction ZIPs but no independently
verified Codabench scores in the preserved record. D222 was built after the
Phase~1 deadline and is therefore excluded from this on-time comparison.

\begin{table}[H]
\centering
\caption{Matched local route variants on separate held-out sets.}
\begin{tabular}{lcc}
\toprule
Variant & Dev BERTScore F1 ($n=144$) & Confirmation F1 ($n=127$) \\
\midrule
Dense control & 0.865228 & 0.865052 \\
Generic prompt route & \textbf{0.868770} & 0.865343 \\
Full role/continuity route & 0.867648 & \textbf{0.869258} \\
\bottomrule
\end{tabular}
\end{table}
The generic route exceeded the full route on the development split, while the
full route scored higher on confirmation. The full route's development
ROUGE-L F1 and BLEU-4 were 0.350857 and 0.264881; on confirmation they were
0.368787 and 0.287066. A paired 5,000-resample bootstrap on archived
per-example BERTScore values gave a full-minus-generic mean difference of
$-0.001122$ on development data (95\% interval $[-0.004886,0.002749]$) and
$0.003915$ on confirmation data (95\% interval
$[-0.001572,0.009617]$). These intervals include zero. The small held-out
sets and selection over checkpoints/decoding policies do not establish a
general causal gain from the full route.

The experiments directly examined an own BPE, four-source sampling,
two training-window types, long-context capacity, a fixed training schedule,
checkpoint/decoding selection, and dense/generic/full routing. They did not
form a full factorial study of tokenizer sizes, source ratios, context
lengths, optimizer settings, or compression. Archived training loss is not a
validation-perplexity curve, and generated-text BERTScore does not verify
legal factuality or citation correctness. The generic route's better local
development score is retained as a negative result for the full-route
hypothesis. The historical model-code and results provenance are preserved in
the October~6 archive and its SHA manifests; the companion notebook in the
code package was assembled later and must not be described as the original
executed GPU notebook.

% Sources to merge into the template bibliography: CS781 Capstone Document
% Phase 1 sections 7.1--7.9; Sennrich et al. 2016 BPE (ACL P16-1162);
% Touvron et al. 2023 LLaMA (arXiv 2302.13971); See et al. 2017
% pointer-generator (ACL P17-1099); Gu et al. 2016 CopyNet (ACL P16-1154).
% Phase 1 code and documentation: https://github.com/k27321278-a11y/cs781-capstone-code
